\subsection{Cuasi-secciones de morfismos planos o lisos.}
\label{subsection:IV.17.16}

Los enunciados de este número completan los de (14.5), mediante hipótesis de planitud.

\begin{proposition}[17.16.1]
\label{IV.17.16.1}
Sea $f:X\to S$ un morfismo plano localmente de presentación finita. Sean $s\in S$, $x$ un punto cerrado de $X_s$ tal que $\sh{O}_{X_s,x}$ sea un anillo de Cohen--Macaulay; entonces existe una vecindad abierta $U$ de $s$ en $S$ y un subpreesquema $X'\subset f^{-1}(U)$ de $X$ tal que $x\in X'$ y que el morfismo $X'\to U$, restricción de $f$, sea plano, cuasi-finito y de presentación finita.
\end{proposition}

La cuestión es local sobre $X$ y $S$, por lo que podemos suponer que $f$ es de presentación finita.\footnote{La transcripción francesa imprime «sur $X$ et $Y$»; el morfismo y el testigo inglés muestran que la base correcta es $S$.} Sea $(t_i)_{1\leq i\leq m}$ un sistema de parámetros del anillo local $\sh{O}_{X_s,x}$; existe una vecindad abierta afín $V=\Spec(A)$ de $x$ en $X$ y $m$ secciones $g_i$ de $\sh{O}_X$ sobre $V$ tales que

\oldpage[IV-4]{106}

las imágenes de los $(g_i)_x\in\sh{O}_{X,x}$ en $\sh{O}_{X_s,x}=\sh{O}_{X,x}/\mathfrak m_s\sh{O}_{X,x}$ sean iguales a los $t_i$. Sea $X'$ el subpreesquema cerrado de $V$ definido por el ideal de $A$ generado por los $g_i$; como la sucesión $(t_i)$ es regular por hipótesis (0, 16.5.7), resulta de (11.3.8) que, reemplazando si es necesario $V$ por una vecindad abierta más pequeña, podemos suponer que el morfismo $X'\to S$, restricción de $f$, es plano y de presentación finita. Por otra parte, puesto que los $t_i$ forman un sistema de parámetros de $\sh{O}_{X_s,x}$, el anillo $\sh{O}_{X'_s,x}$ es artiniano por definición, y como $x$ es cerrado en $X'_s$, se concluye que está aislado en $X'_s$. Reemplazando $V$ si es necesario por una vecindad más pequeña de $x$, se concluye, gracias a (13.1.4), que el morfismo $X'\to S$ es cuasi-finito.

\begin{corollary}[17.16.2]
\label{IV.17.16.2}
Sea $f:X\to S$ un morfismo fielmente plano y localmente de presentación finita. Entonces existe un morfismo $g:S'\to S$, fielmente plano, localmente de presentación finita y localmente cuasi-finito, tal que exista un $S$-morfismo $S'\to X$ (en otras palabras, que exista una $S'$-sección de $X'=X\times_S S'$ (I, 3.3.14)). Si $S$ es cuasi-compacto (resp. cuasi-compacto y cuasi-separado), se puede suponer que $S'$ es afín (resp. $S'$ afín y el morfismo $g$ de presentación finita y cuasi-finito).\footnote{La frase francesa omite «de présentation finie»; se restituye según el testigo inglés y la condición de la proposición.}
\end{corollary}

Para todo $s\in S$, la fibra $X_s$ no es vacía por hipótesis y es un preesquema localmente de tipo finito sobre $k(s)$; el conjunto $U_s$ de los puntos $x$ de $X_s$ donde $\sh{O}_{X_s,x}$ es un anillo de Cohen--Macaulay es abierto en $X_s$ (6.11.3) y no es vacío, puesto que contiene los puntos maximales de $X_s$ (0, 16.5.1); contiene, por tanto, un punto $x_s$ cerrado en $X_s$ (10.4.7). Sea $X'(s)$ un subpreesquema afín de $X$ que contenga $x_s$ y verifique las condiciones de (17.16.1). Para obtener un preesquema $S'$ que verifique las condiciones del enunciado, basta tomar la suma de los $X'(s)$, recorriendo $s$ el conjunto $S$. Puesto que el morfismo $X'(s)\to S$ es plano y localmente de presentación finita, la imagen $U(s)$ de $X'(s)$ es abierta en $S$ (2.4.6); cuando $S$ es cuasi-compacto, podemos, pues, recubrir $S$ mediante un número finito de los $U(s_i)$ y el preesquema $S'$ suma de los $X'(s_i)$ sigue respondiendo a la cuestión y es afín. Si además $S$ es cuasi-separado, podemos suponer que las inmersiones abiertas $U(s_i)\to S$ son cuasi-compactas (I, 2.7), por tanto de presentación finita (I, 6.2), y entonces los morfismos $X'(s_i)\to S$ son de presentación finita (I, 6.2), y también lo es, por consiguiente, el morfismo $S'\to S$.

\begin{corollary}[17.16.3]
\label{IV.17.16.3}
\begin{enumerate}
  \item[(i)] Sea $f:X\to S$ un morfismo liso. Sean $s\in S$, $x$ un punto cerrado de $X_s$ tal que el cuerpo residual $k(x)$ sea separable sobre $k(s)$; entonces, en la conclusión de (17.16.1), se puede tomar $X'$ de modo que el morfismo $X'\to U$, restricción de $f$, sea étale.
  \item[(ii)] Sea $f:X\to S$ un morfismo liso sobreyectivo. Entonces, en las conclusiones de (17.16.2), se puede suponer además que $g:S'\to S$ sea étale.
\end{enumerate}
\end{corollary}

Es claro que (ii) se deduce de (i) del mismo modo que (17.16.2) de (17.16.1), teniendo en cuenta que, para todo $s\in S$, como $X_s$ no es vacío y es liso sobre $k(s)$, existe un punto cerrado $x\in X_s$ tal que $k(x)$ sea separable sobre $k(s)$ (17.15.10, (iii)). Basta, pues, demostrar (i). Se observará que el anillo $\sh{O}_{X_s,x}$ es entonces regular; si se repite la construcción hecha en (17.16.1) tomando para $(t_i)$ un sistema regular de parámetros de $\sh{O}_{X_s,x}$, el anillo $\sh{O}_{X'_s,x}$ es un cuerpo isomorfo a $k(x)$, por tanto separable sobre $k(s)$ por hipótesis. La conclusión resulta entonces de (17.6.1, $c'$)).

\begin{proposition}[17.16.4]
\label{IV.17.16.4}
Sean $S$ un preesquema cuasi-compacto y cuasi-separado, $f:X\to S$ un morfismo sobreyectivo localmente de presentación finita. Entonces existe una familia finita $(S_i)_{i\in I}$

\oldpage[IV-4]{107}

de subpreesquemas afines de $S$, de presentación finita sobre $S$, dos a dos disjuntos, cuya unión es $S$, y que tiene la propiedad siguiente: para cada $i\in I$, existen dos morfismos finitos, de presentación finita y sobreyectivos
\[
  S_i''\xrightarrow{g_i}S_i'\xrightarrow{h_i}S_i,
\]
donde $h_i$ es étale y $g_i$ es plano y radicial, y un $S$-morfismo $S_i''\to X$ (en otras palabras, una $S_i''$-sección de $X_i''=X\times_S S_i''$).
\end{proposition}

Hay un recubrimiento finito $(U_i)_{1\leq i\leq n}$ de $S$ por abiertos afines; para todo $i$, la intersección
\[
  W_i=U_i\cap\left(\bigcup_{j<i}U_j\right)
\]
es cuasi-compacta (I, 2.7), por lo que existe un subpreesquema cerrado $V_i$ de $U_i$ cuyo espacio subyacente es $U_i-W_i$ y que está definido por un ideal de tipo finito del anillo de $U_i$; por tanto, $V_i$ es un subpreesquema afín de $S$ que es de presentación finita sobre $S$ (I, 6.2). Es claro que podemos limitarnos a demostrar la proposición reemplazando $S$ por $V_i$ y $X$ por $X\times_S V_i$. En otras palabras, ya podemos suponer que $S$ es afín. Por otra parte, $X$ es unión de abiertos afines $W_\alpha$ y los $f(W_\alpha)$ son constructibles en $S$ (I, 8.4) y forman un recubrimiento de $S$; por consiguiente (I, 9.9), existe una subfamilia finita $(W_{\alpha_j})_{1\leq j\leq m}$ tal que los $f(W_{\alpha_j})$ ya forman un recubrimiento de $S$. Sea entonces $X'$ el preesquema suma de los subpreesquemas inducidos sobre los abiertos $W_{\alpha_j}$ de $X$; es inmediato que basta demostrar la proposición reemplazando $X$ por $X'$, puesto que un $S$-morfismo $S_i''\to X'$ da por composición un $S$-morfismo $S_i''\to X'\to X$.

Podemos, pues, suponer $X$ afín y $f$ de presentación finita. Se puede entonces (8.9.1) escribir $X$ bajo la forma $X_0\times_{S_0}S$, donde $S_0$ es noetheriano, $X_0\to S_0$ un morfismo de tipo finito, que podemos suponer sobreyectivo (8.10.5, (vi)). Si se prueba la proposición para este morfismo, resultará inmediatamente para $X$ por cambio de base $S\to S_0$. Podemos, pues, suponer además que $S$ es noetheriano y que $f$ es de tipo finito.

Sean $s_h$ $(1\leq h\leq r)$ los puntos maximales de $S$. Demostremos que basta probar el enunciado reemplazando $S$ por una vecindad abierta suficientemente pequeña $U_h$ de cada uno de los $s_h$, y $X$ por la imagen inversa de $U_h$ en $X$. En efecto, supongamos establecida la proposición en este caso, y razonemos por recurrencia noetheriana ($0_{\mathrm I}$, 2.2.2), suponiendo establecido el enunciado para todo subpreesquema cerrado de $S$ cuyo espacio subyacente sea $\neq S$. Podemos suponer que los $U_h$ son dos a dos disjuntos. Si $T$ es un subpreesquema cerrado cuyo espacio subyacente es $S-(\bigcup_h U_h)$, la hipótesis de recurrencia entraña que el enunciado es cierto para $T$; como también es cierto para cada uno de los $U_h$, es evidentemente cierto para $S$.

Como evidentemente podemos (reemplazando $S$ por $S_{\mathrm{red}}$ y $X$ por $X\times_S S_{\mathrm{red}}$) suponer que $S$ es reducido, cada uno de los $\sh{O}_{S,s_h}$ es un cuerpo $k(s_h)$. Supongamos que se ha demostrado la proposición cuando $S$ es el espectro de un cuerpo y $f$ es de tipo finito. Entonces la existencia de los $U_h$ resulta del método de (8.1.2, $a$)) y de (8.8.2, (i) y (ii)), (8.10.5, (vi), (vii) y (x)), (11.2.6, (ii)) y (17.7.8, (ii)).

Supongamos, pues, $S=\Spec(k)$, donde $k$ es un cuerpo, siendo $X$ de tipo finito sobre $k$ y $X\neq\varnothing$. Como $X$ es noetheriano, existe en $X$ un punto cerrado $x$ ($0_{\mathrm I}$, 2.1.3), por lo que $k(x)$ es una extensión finita de $k$ (I, 6.4.2). Existe, por tanto, una extensión finita separable $k'$ de $k$ tal que $k''=k(x)$ sea una extensión finita radicial de $k'$. Respondemos entonces a la cuestión tomando $I$ reducido a un elemento, $S_i'=\Spec(k')$, $S_i''=\Spec(k'')$.

\oldpage[IV-4]{108}

\begin{corollary}[17.16.5]
\label{IV.17.16.5}
Sea $f:X\to S$ un morfismo sobreyectivo localmente de presentación finita. Entonces existe un morfismo $g:S'\to S$ sobreyectivo, localmente de presentación finita y localmente cuasi-finito, tal que exista un $S$-morfismo $S'\to X$ (en otras palabras, que exista una $S'$-sección de $X'=X\times_S S'$). Si $S$ es cuasi-compacto (resp. cuasi-compacto y cuasi-separado), se puede suponer que $S'$ es afín (resp. $S'$ afín y el morfismo $g$ de presentación finita y cuasi-finito).
\end{corollary}

Es suficiente demostrar el corolario suponiendo que $S$ es afín: en efecto, $S$ es unión de una familia $(S_\alpha)$ de abiertos afines, y si para cada $\alpha$, $S'_\alpha$ es afín y el morfismo $g_\alpha:S'_\alpha\to S_\alpha$ responde a la cuestión y es de presentación finita y cuasi-finito, tomando para $S'$ el preesquema suma de los $S'_\alpha$, con $g:S'\to S$ coincidiendo con $g_\alpha$ en cada uno de los $S'_\alpha$, este morfismo responde a la cuestión; además, si $S$ es cuasi-compacto, podemos suponer finita la familia $(S_\alpha)$, por lo que $S'$ es afín; si además $S$ es cuasi-separado, las inmersiones $S_\alpha\to S$ son de presentación finita (I, 6.2), por lo que también lo es $g$.

Consideremos, pues, el caso en que $S$ es afín: formamos entonces los morfismos finitos y de presentación finita $S_i''\to S_i$ de (17.16.4); como las inmersiones $S_i\to S$ son de presentación finita (siendo los $S_i$ afines), los morfismos compuestos $g_i'':S_i''\to S_i\to S$ son de presentación finita y cuasi-finitos, y respondemos a la cuestión tomando $S'$ igual a la suma de los $S_i''$ y $g$ coincidiendo con $g_i''$ en cada uno de los $S_i''$.

\begin{proposition}[17.16.6]
\label{IV.17.16.6}
Sean $S$ un preesquema cuasi-compacto y cuasi-separado, $f:X\to S$ un morfismo de presentación finita tal que, para todo $s\in S$, $X_s=f^{-1}(s)$ sea un preesquema propio sobre $k(s)$. Entonces existe una familia finita $(S_i)_{i\in I}$ de subpreesquemas afines de $S$, de presentación finita sobre $S$, dos a dos disjuntos, cuya unión es $S$, y tal que, para todo $i$, el morfismo $X_i=X\times_S S_i\to S_i$ sea propio y plano. Si además, para todo $s\in S$, $X_s$ es finito sobre $k(s)$, se pueden tomar los $S_i$ de modo que cada uno de los morfismos $X_i\to S_i$ se factorice como
\[
  X_i\xrightarrow{u_i}X_i'\xrightarrow{h_i}S_i
\]
donde $u_i$ y $h_i$ son finitos y localmente libres (18.2.7), $h_i$ es étale, y $u_i$ es radicial y sobreyectivo.
\end{proposition}

Seguimos un procedimiento análogo al de (17.16.4). Nos reducimos primero al caso en que $S$ es afín y $f$ es plano, utilizando el teorema de planitud genérica (8.9.5) (obsérvese que los subpreesquemas $S_i'$ definidos en la demostración de (8.9.5) son afines). Puesto que $f$ es de presentación finita, podemos escribir a continuación $X=X_0\times_{S_0}S$, donde $S_0$ es noetheriano y $X_0\to S_0$ un morfismo de tipo finito y plano (11.2.7); además, cada fibra $(X_0)_{s_0}$ sigue siendo propia sobre $k(s_0)$, como resulta de (2.7.1, (vii)), y nos reducimos, por tanto, al caso en que $S$ es noetheriano. Por recurrencia noetheriana y aplicación del procedimiento de (8.1.2, $a$)) (utilizando también esta vez (8.10.5, (xii))), para demostrar la primera afirmación nos reducimos finalmente al caso en que $S$ es el espectro de un cuerpo, que resulta trivialmente de la hipótesis (con $S_i=S$). Se trata de la misma manera el caso en que se supone que $X_s$ es finito sobre $k(s)$ para todo $s\in S$ (esta vez hay que utilizar (2.7.1, (xv)), (8.10.5, (x), (iv), (vi) y (vii)), (17.7.8) y (2.1.12)), y nos reducimos a demostrar la última afirmación del enunciado cuando $S=\Spec(k)$ es el espectro de un cuerpo.

Pero entonces $X=\Spec(A)$, donde $A$ es una $k$-álgebra finita (I, 6.4.4); por tanto, $A$ es un producto directo de $k$-álgebras finitas $A_j$ que son anillos locales $(1\leq j\leq m)$;\footnote{La transcripción francesa dice «composée directe»; aquí se expresa el producto directo indicado por el testigo inglés y por la descomposición de una álgebra artiniana.} puesto que $A_j$ es un anillo artiniano y una

\oldpage[IV-4]{109}

$k$-álgebra, contiene un subcuerpo $K'_j$ isomorfo canónicamente al cuerpo residual de $A_j$ (0, 19.6.2). Tomamos entonces $X_i=X$, $X'_i$ suma de los $\Spec(K'_j)$ y es claro que respondemos así a la cuestión, pues para todo $j$ tenemos dos homomorfismos $K'_j\to A_j\to k(A_j)$ cuyo compuesto es un isomorfismo.
